%% LyX 2.5.1 created this file.  For more info, see https://www.lyx.org/.
%% Do not edit unless you really know what you are doing.
\documentclass[english,hebrew]{article}
\usepackage[HE8,T1]{fontenc}
\usepackage[utf8]{inputenc}
\usepackage{babel}
\usepackage{amsmath}
\usepackage{amssymb}
\usepackage[unicode=true]
 {hyperref}
\begin{document}
\title{מבוא לטופולוגיה קבוצתית: קשירות}
\maketitle
\begin{description}
\item [{קטגוריות:}] טופולוגיה
\item [{תגים:}] קשירות
\item [{מזהה:}] \L{connectedness}
\end{description}

\section{מה זה בכלל?}

סדרת הפוסטים שלי על טופולוגיה קבוצתית \href{https://gadial.net/2026/07/28/topological_spaces_compactness/}{הגיעה אל מושג הקומפקטיות}
ויצאה שם להפסקה. בואו נחזור עם נושא אחר שגם הוא באותה רוח, אבל קצת
יותר פשוט: \textbf{קשירות} )\L{connectedness}( של מרחב טופולוגי. כמו
עם קומפקטיות, גם פה אפשר לתת מוטיבציה ראשונית לנושא מתוך משפט בסיסי
על פונקציות רציפות שרואים בחדו\textquotedblright א; אם עבור קומפקטיות
אלו היו משפטי ויירשטראס, הפעם זה יהיה \textbf{משפט ערך הביניים}. בואו
נזכיר אותו:

משפט ערך הביניים בחדו\textquotedblright א אומר שאם יש לנו פונקציה
\L{$f:\left[a,b\right]\to\mathbb{R}$} שהיא רציפה )כלומר - היא רציפה
על הקטע הסגור \L{$\left[a,b\right]$}; היא יכולה להיות מוגדרת גם מחוץ
לקטע הזה אבל לא חשוב לנו אם היא רציפה או לא שם( אז היא מחזירה על \L{$\left[a,b\right]$}
כל ערך שהוא בין \L{$f\left(a\right)$} ו-\L{$f\left(b\right)$}, כלומר
כל ערך שהוא \textquotedblleft בין לבין\textquotedblright{} שני הערכים
הללו, ומכאן שם המשפט. פורמלית, לכל \L{$y\in\mathbb{R}$} שמקיים \L{$f\left(a\right)\le y\le f\left(b\right)$}
או \L{$f\left(b\right)\le y\le f\left(a\right)$} )זה תלוי בשאלה מי
גדול ממי מבין \L{$f\left(a\right),f\left(b\right)$}( קיים \L{$c\in\left[a,b\right]$}
כך ש-\L{$f\left(c\right)=y$}.

זה משפט שימושי וחשוב מאוד; שימוש לא טריוויאלי חמוד מאוד שלו הוא להראות
שאם \L{$p\left(x\right)$} הוא פולינום ממעלה אי זוגית, קיים לו שורש
ממשי, כלומר \L{$c$} כך ש-\L{$p\left(c\right)=0$}. זה עובד בגלל שפולינום
הוא פונקציה רציפה, וגם כשמשאיפים את \L{$x$} לאינסוף וכשמשאיפים אותו
למינוס אינסוף, אז החזקה הגבוהה ביותר של הפולינום תשאף פעם אחת לאינסוף
ופעם אחת למינוס אינסוף )להבדיל מהמקרה של פולינום ממעלה זוגית, שבה
שני הכיוונים ישאפו לאותו דבר, זה רק תלוי בשאלה אם המקדם המוביל הוא
חיובי או שלילי(. בפרט אפשר למצוא \L{$a,b$} כך ש-\L{$p\left(a\right)\le0\le p\left(b\right)$}
ומשפט ערך הביניים משלים את העבודה. \href{https://gadial.net/2009/10/29/fundemental_theorem_of_algebra_algebraic_proof/}{בפוסט שלי על ההוכחה האלגברית}
למשפט היסודי של האלגברה אני מציג בדיוק את השימוש הזה בתור הצעד הראשון
בהוכחה )הצעד ה\textquotedblright לא אלגברי\textquotedblright{} שחייבים
לבצע לפני שמגיעים להוכחה האלגברית עצמה(.

איך מוכיחים את המשפט? ראשית, שמים לב שהוא לא נכון אם התחום של הפונקציה
הוא המספרים הרציונליים \L{$\mathbb{Q}$}. במקרה הזה אפשר לקחת פונקציה
קבועה שהרציפות שלה \textquotedblleft נשברת\textquotedblright{} והיא
מבצעת קפיצה על נקודה לא רציונלית ומכיוון שהפונקציה מוגדרת רק על הרציונליים
היא לא תשים לב לזה. הנה דוגמא. נגדיר \L{$f:\mathbb{Q}\to\mathbb{Q}$}
על ידי

\L{$f\left(x\right)=\begin{cases}
0 & x^{2}<2\\
1 & x^{2}\ge2
\end{cases}$}

כלומר, הפונקציה \textquotedblleft שוברת את הרציפות שלה\textquotedblright{}
על \L{$\sqrt{2}$}, אבל כל עוד העולם שלנו הוא \L{$\mathbb{Q}$}, הנקודה
\L{$\sqrt{2}$} לא קיימת בכלל; לעומת זאת הקריטריון \L{$x^{2}<2$}
הוא דרך קונקרטית לבדוק ש-\L{$x<\sqrt{2}$} גם בלי לדעת שהנקודה הזו
קיימת. עכשיו, קחו כל נקודה \L{$x_{0}\in\mathbb{Q}$} שהיא; אם \L{$x^{2}_{0}<2$}
אז קיימת לה סביבה שכל נקודה בה עדיין מקיימת \L{$x^{2}<2$} ולכן הפונקציה
תהיה רציונלית על \L{$x_{0}$}, וכך גם אם \L{$x^{2}_{0}>2$}. המקרה
הבעייתי האפשרי היחיד הוא \L{$x^{2}_{0}=2$} שכאמור, לא יכול לקרות.
אז קיבלנו פונקציה רציפה על הרציונליים שמקבלת {\beginL 0\endL} ומקבלת
{\beginL 1\endL} אבל לא מקבלת \L{$\frac{1}{2}$} בשום מקום, בסתירה
למשפט.

מה השתבש כאן? בבירור יש לנו בעיה עם \textbf{התחום} של הפונקציה. הצלחנו
איכשהו לחלק אותו לשני חלקים שונים שנקודת ההפרדה ביניהם היא בכלל לא
חלק מהמרחב ולכן אנחנו לא רואים אותה. חילקנו את המרחב לשתי \textbf{קבוצות
פתוחות} זרות: \L{$\left(-\infty,\sqrt{2}\right)$} ו-\L{$\left(\sqrt{2},\infty\right)$}.
הקטע הזה של חלוקה של המרחב לשתי קבוצות פתוחות זרות? זו הולכת להיות
\textbf{ההגדרה} שנשתמש בה בהמשך.

למה ב-\L{$\mathbb{R}$} זה לא קורה? אינטואיטיבית, אנחנו אומרים שב-\L{$\mathbb{R}$}
\textquotedblleft אין חורים\textquotedblright{} בגלל שהוא \textquotedblleft רציף\textquotedblright .
פורמלית אנחנו מדברים על מה שנקרא \textbf{אקסיומת השלמות} שאני מזכיר
\href{https://gadial.net/2024/10/12/real_numbers_two_completness/}{בפוסט הזה}
וכבר דיברנו עליה כשדיברנו על קומפקטיות. האקסיומה הזו לא מאפשרת למשהו
כמו \L{$\sqrt{2}$} פשוט לא להתקיים: אם אני יכול לדבר על קבוצה כמו
\L{$\left\{ x\in\mathbb{R}|x^{2}<2\right\} $} אז גם החסם העליון שלה,
\L{$\sqrt{2}$}, הולך להיות קיים. ככה עובדת בערך ההוכחה של משפט ערך
הביניים: מסתכלים על הקבוצה \L{$\left\{ x\in\left[a,b\right]|f\left(x\right)\le y\right\} $},
מראים שיש לה חסם עליון \L{$c$}, ומשתמשים ברציפות של \L{$f$} כדי
להראות ש-\L{$f\left(c\right)=y$}.

עכשיו אנחנו רוצים לקחת את כל זה ולהכליל למרחבים טופולוגיים. במרחב
טופולוגי לא תהיה לנו את אקסיומת השלמות, אבל כל הרעיון הוא לא להתבסס
עליה כי היא תכונה \textquotedblleft חדו\textquotedblright אתית\textquotedblright{}
של הישר הממשי, אלא לזהות מה \textbf{התכונה הטופולוגית} הרלוונטית;
וכבר ראינו אותה, זו התכונה \textquotedblleft אי אפשר לפרק את המרחב
שלנו לשתי קבוצות פתוחות זרות\textquotedblright , אז בואו נגדיר פורמלית:
\begin{itemize}
\item מרחב טופולוגי \L{$X$} נקרא \textbf{קשיר} אם לא קייימות \L{$A,B\subseteq X$}
שהן קבוצות פתוחות זרות )\L{$A\cap B=\emptyset$}( שאיחודן כל \L{$X$}
)\L{$A\cup B=X$}(.
\end{itemize}
יש לזה קריטריון שקול חמוד. כזכור, \textbf{קבוצה סגורה} במרחב טופולוגי
היא קבוצה שהמשלימה שלה היא קבוצה פתוחה )ומה זו קבוצה פתוחה? כמובן,
קבוצה שהמשלימה שלה סגורה(. שתי הקבוצות הפתוחות שיש בכל מרחב טופולוגי
\L{$X$} הן \L{$\emptyset$} ו-\L{$X$} עצמה והן המשלימות זו של זו,
כך ששתיהן פתוחות וסגורות בו זמנית, מה שבאנגלית מכונה בכינוי \L{clopen}
האלגנטי ובעברית מכונה \textquotedblleft פסגורה\textquotedblright{}
שזה ביטוי שלא כדי להגיד בהרצאה כדי לא להתיז רוק על השורות הראשונות
וסלחו לי אבל לא אשתמש לא בזה ולא בזה.

האם יש עוד קבוצות שהן פתוחות וסגורות בו זמנית? אה-הא! זה שקול לטענה
ש-\L{$X$} \textbf{לא קשיר}. אם \L{$X$} לא קשיר אז \L{$A,B$} הפתוחות
הללו הן משלימות זו של זו )תוכיחו!( ולכן שתיהן גם סגורות, כך שמרחב
לא קשיר מכיל קבוצות כאלו. ואם יש לנו מרחב \L{$X$} שמכיל קבוצה \L{$A$}
שהיא גם פתוחה וגם סגורה, אז נגדיר \L{$B=X\backslash A$} ונקבל מכך
ש-\L{$A$} סגורה ש-\L{$B$} פתוחה וקיבלנו פירוק של המרחב לאיחוד זר
של שתי קבוצות פתוחות ולכן המרחב לא קשיר. זה היה פשוט אבל חמוד.

\section{משפט ערך הביניים}

עכשיו, איך ננסח גרסה מוכללת של משפט ערך הביניים? אנחנו לוקחים פונקציה
\L{$f:X\to Y$} שהיא רציפה כך ש-\L{$X$} קשיר; ואנחנו לוקחים שתי נקודות
\L{$a,b\in X$}, ועכשיו אנחנו רוצים נקודת \textbf{ביניים} \L{$y\in Y$}
שנמצאת בין \L{$f\left(a\right)$} ו-\L{$f\left(b\right)$}. בשביל
בכלל לדבר על נקודה כזו, אנחנו צריכים שעל \L{$Y$} יהיה מבנה של \textbf{סדר}
כלשהו, ואם אנחנו רוצים הוכחה טופולוגית, הטופולוגיה של \L{$Y$} צריכה
איכשהו להתקשר אל הסדר שמוגדר על \L{$Y$}. למרבה המזל יש לנו בדיוק
את הדבר הזה - \textbf{טופולוגיית הסדר} שכבר \href{https://gadial.net/2026/07/07/topological_spaces_and_bases/}{דיברתי עליה כאן}.
אז הנה הניסוח הפורמלי של המשפט שלנו:
\begin{itemize}
\item )משפט ערך הביניים(: אם \L{$f:X\to Y$} היא פונקציה רציפה כך ש-\L{$X$}
מרחב טופולוגי קשיר ו-\L{$Y$} מרחב טופולוגי סדור עם טופולוגיית הסדר,
ואם יש \L{$a,b\in X$} ו-\L{$y\in Y$} כך ש-\L{$f\left(a\right)\le y\le f\left(b\right)$},
אז קיים \L{$c\in X$} כך ש-\L{$f\left(c\right)=y$}.
\end{itemize}
איך מוכיחים את זה? בחדו\textquotedblright א ההוכחה הסתמכה על כל מני
תעלולי אפסילון-דלתא מעיקים של פונקציות רציפות, אבל אצלנו זה פשוט לא
קיים. מראש דרך החשיבה צריכה להיות שונה, \textquotedblleft גלובלית\textquotedblright{}
יותר. אז הנה שני הרעיונות שיתנו לנו את המשפט:
\begin{itemize}
\item תמונה של מרחב קשיר על ידי פונקציה רציפה הוא מרחב קשיר, כלומר \L{$f\left(X\right)\subseteq Y$}
יהיה קשיר.
\item תת-מרחב קשיר בטופולוגיית הסדר לא מכיל \textquotedblleft חורים באמצע\textquotedblright .
\end{itemize}
את הטענה השנייה קל לראות מייד: אם \L{$Y$} סדור בטופולוגיית הסדר,
ו-\L{$W\subseteq Y$} הוא תת-מרחב קשיר, בואו ניקח \L{$a\in Y$} כלשהו
ונסתכל על שתי הקבוצות הבאות של כל מה שגדול וכל מה שקטן מ-\L{$a$}:

\L{$\left(-\infty,a\right)=\left\{ y\in Y\ |\ y<a\right\} $}

\L{$\left(a,\infty\right)=\left\{ y\in Y\ |\ a<y\right\} $}

בטופולוגיית הסדר, שתי הקבוצות הללו הן קבוצות פתוחות. עכשיו בואו נסתכל
על החיתוך שלהן עם \L{$W$}, כלומר על הקבוצות

\L{$W\cap\left(-\infty,a\right)=\left\{ y\in W\ |\ y<a\right\} $}

\L{$W\cap\left(a,\infty\right)=\left\{ y\in W\ |\ a<y\right\} $}

הטופולוגיה של \L{$W$} היא מה שנקרא \textbf{טופולוגיית תת-המרחב} שתיארתי
\href{https://gadial.net/2026/07/25/topological_spaces_subspace_and_product_topology/}{בפוסט הזה},
ובטופולוגיה הזו הקבוצות הפתוחות הן בדיוק חיתוך של קבוצות פתוחות במרחב
המקורי עם תת-המרחב, כלומר שתי הקבוצות שתיארתי עכשיו פתוחות ב-\L{$W$}.
הן גם זרות, כי אין \L{$y$} כך ש-\L{$y<a$} וגם \L{$a<y$}. אחת מהן
\textbf{יכולה} להיות ריקה; אם למשל אין ב-\L{$W$} איברים קטנים מ-\L{$a$}
הקבוצה הראשונה תהיה ריקה, אבל במקרה הזה \L{$a$} הוא לא \textquotedblleft באמצע\textquotedblright{}
\L{$W$} אלא בהתחלה )או לפני ההתחלה( של \L{$W$}. כמו כן, האיחוד של
הקבוצות לא בהכרח נותן את כל \L{$W$}; אם \L{$a\in W$} אז הוא יישאר
בחוץ. עכשיו, אם \L{$W$} קשיר זה אומר ש\textbf{לא ייתכן} ששתי הקבוצות
לא ריקות והאיחוד שלהן נותן את כל \L{$W$} כי אז קיבלנו פירוק של \L{$W$};
לכן, אם יש \L{$a\notin W$} כלשהו, בהכרח \L{$a<x$} לכל \L{$x\in W$}
או \L{$a>x$} לכל \L{$x\in W$}. זו המשמעות של \textquotedblleft אין
חורים באמצע\textquotedblright .

עכשיו נחזור להוכחה המקורית: \L{$X$} קשיר, \L{$f$} רציפה, נסמן \L{$f\left(X\right)=W$},
נשתמש בטענה שטרם הוכחתי לפיה נובע ש-\L{$W$} קשיר, ניקח \L{$a,b\in X$},
ו-\L{$y\in Y$} כך ש-\L{$f\left(a\right)\le y\le f\left(b\right)$}
ונסיק מהקשירות של \L{$W$} שבהכרח \L{$y\in W$} כי אם היה מתקיים \L{$y\notin W$}
הוא לא היה יכול בו זמנית להיות גם גדול וגם קטן מאיברים ב-\L{$W$}.
מכיוון ש-\L{$y\in W=f\left(X\right)$} קיים \L{$c\in X$} כך ש-\L{$f\left(c\right)=y$}
וסיימנו.

יפה, אז כדי להשלים את ההוכחה של משפט ערך הביניים המוכלל צריך להוכיח
שתמונה רציפה של מרחב קשיר הוא קשיר. שימו לב: לא תמונה הומיאומורפית.
אם שני מרחבים הם הומיאומורפיים, כלומר קיימת בינם פונקציה חח\textquotedblright ע
ועל שהיא וההופכית שלה שתיהן רציפות, זה אומר שבשני המרחבים יש בדיוק
את אותן קבוצות פתוחות ולכן ברור שאם אחד קשיר גם השני קשיר. כאן אנחנו
צריכים להתבסס על פחות - הפונקציה ההפוכה לא חייבת להיות רציפה או קיימת
בכלל כי ייתכן שהפונקציה שלנו לא חח\textquotedblright ע. עדיין, האינטואיציה
ברורה - הרציפות לא מאפשרת \textquotedblleft לקרוע\textquotedblright{}
מרחב שהוא חתיכה אחת לשתי חתיכות שונות.

בואו נניח ש-\L{$X$} קשיר ו-\L{$f$} פונקציה רציפה מעל \L{$X$} ונסמן
\L{$f\left(X\right)=Y$}. אני רוצה להוכיח ש-\L{$Y$} קשיר, אז בואו
נסתכל על פירוק \L{$Y=A\cup B$} כך ש-\L{$A\cap B=\emptyset$} ושתי
הקבוצות הללו פתוחות, ונוכיח שלפחות אחת מהן ריקה.

כל מה שיש לנו להסתמך עליו הוא ש-\L{$f$} רציפה. מה זה אומר? זה אומר
ש-\L{$f^{-1}\left(A\right)=\left\{ x\in X\ |\ f\left(x\right)\in A\right\} $}
היא קבוצה פתוחה ובדומה גם \L{$f^{-1}\left(B\right)$} פתוחה. עכשיו,
\href{https://gadial.net/2026/07/12/topological_spaces_closure_and_continuous_functions/}{בפוסט שלי על פונקציות רציפות}
כבר ראינו שמתקיים

\L{$f^{-1}\left(A\right)\cup f^{-1}\left(B\right)=f^{-1}\left(A\cup B\right)=$}

\L{$=f^{-1}\left(Y\right)=X$}

כלומר, קיבלנו שתי קבוצות פתוחות שהאיחוד שלהן הוא \L{$X$}. האם הן
זרות? בוודאי, כי 

\L{$f^{-1}\left(A\right)\cap f^{-1}\left(B\right)=f^{-1}\left(A\cap B\right)=f^{-1}\left(\emptyset\right)=\emptyset$}

אם כן, בגלל ש-\L{$X$} קשיר, אחת משתי הקבוצות הללו חייבת להיות ריקה,
נאמר \L{$f^{-1}\left(A\right)=\emptyset$}. אבל מכיוון ש-\L{$A\subseteq f\left(X\right)$},
\textbf{לכל} \L{$y\in A$} קיים \L{$x\in X$} כך ש-\L{$f\left(x\right)=y$},
ולכן אם \L{$f^{-1}\left(A\right)=\emptyset$} אז \L{$A=\emptyset$},
מה שמשלים את ההוכחה.

זה מסיים עם המוטיבציה לפוסט הזה - הכללת משפט ערך הביניים. אבל אם אנחנו
כבר כאן, יש עוד כמה דברים בסיסיים מועילים לומר על קשירות.

\section{בונים קבוצות קשירות מקבוצות קשירות}

ראשית, בטופולוגיה אוהבים לבנות מרחבים חדשים מתוך מרחבים קיימים, ואם
הרכיבים שלנו קשירים זה תמיד מעניין לראות האם גם התוצאה תהיה קשירה.
דרך אחת לבנות מרחב חדש ממרחב קיים כבר ראינו - להפעיל עליו פונקציה
רציפה, מה שאכן משמר קשירות. מה עם איחוד של כמה תתי-מרחבים? ברור שבאופן
כללי הוא לא ישמר קשירות - הקטעים \L{$\left(0,1\right)$} ו-\L{$\left(2,3\right)$}
בישר הממשי הם קשירים כל אחד בנפרד )בעצם... עוד לא הוכחתי את זה! אבל
נחכה עם זה להמשך( ואילו \L{$\left(0,1\right)\cup\left(2,3\right)$}
הוא בבירור לא קשיר כי הנה, זה עתה ראינו פירוק שלו. מה שכן, אם הייתה
להם נקודה משותפת, אז כבר ברור שהאיחוד היה יוצא קשיר. וזה נכון באופן
כללי, לכל מספר של תת-מרחבים שאנחנו מאחדים:
\begin{itemize}
\item אם \L{$X$} מרחב טופולוגי ו-\L{$\left\{ W_{i}\right\} _{i\in\Lambda}$}
היא משפחה כלשהי של תת-מרחבים קשירים של \L{$W$} כך ש-\L{$\bigcap_{i\in\Lambda}W_{i}\ne\emptyset$}
אז \L{$\bigcup_{i\in\Lambda}W_{i}$} קשיר.
\end{itemize}
זה די מפתיע, אותי לפחות, כי איכשהו התחושה האינטואיטיבית שלי היא שקשירות
היא דווקא לא תכונה \textquotedblleft נקודתית\textquotedblright{} כי,
נו, אנחנו מדברים על פירוק של המרחב לקבוצות פתוחות והן שמנות כאלה,
נקודה בודדת זה יותר קטע של קבוצות סגורות. כמובן שזו אינטואיציה רעה
וכבר ראינו קודם שקשירות ניתנת לניסוח גם עם שתי קבוצות סגורות ולא פתוחות,
ועוד מעט נראה עוד סוג של קשירות שגם בו עניינים \textquotedblleft נקודתיים\textquotedblright{}
דווקא מסתדרים טוב, אבל בואו נעבור עכשיו להוכחה של הטענה עצמה.

נניח שיש ל-\L{$\bigcup_{i\in\Lambda}W_{i}$} פירוק לשתי קבוצות פתוחות
זרות \L{$A,B$}. עבור \L{$W_{k}$} ספציפי, הקבוצות \L{$A_{k}=W_{k}\cap A$}
ו-\L{$B_{k}=W_{k}\cap B$} הן פתוחות ב-\L{$W_{k}$} לפי טופולוגיית
תת-המרחב. כמו כן \L{$W_{k}\subseteq\bigcup_{i\in\Lambda}W_{i}=A\cup B$}
ולכן

\L{$A_{k}\cup B_{k}=W_{k}\cap\left(A\cup B\right)=W_{k}$}

אז קיבלנו שתי קבוצות פתוחות זרות שאיחודן הוא המרחב הקשיר \L{$W_{k}$}
ולכן \L{$A_{k}=\emptyset$} או \L{$B_{k}=\emptyset$}. בואו נניח בלי
הגבלת הכלליות ש-\L{$A_{k}=\emptyset$}, אז באותו אופן גם \L{$A_{i}=\emptyset$}
לכל \L{$i\in\Lambda$} ולכן \L{$A=\bigcup A_{i}=\emptyset$}... רגע
רגע רגע, זה לא נכון. זה שהנחתי בלי הגבלת הכלליות ש-\L{$A_{k}=\emptyset$}
לא אומר שאני יכול להניח את זה \textbf{לכל} \L{$i\in\Lambda$}; זה
בדיוק \textbf{ויתור} על המקרה הכללי, שבו ייתכן שבחלק מהקבוצות זה יהיה
\L{$A_{i}$} ובחלק מהן זה יהיה \L{$B_{i}$}.

למרבה המזל, קל להראות שזה לא באמת קורה, בעזרת הנתון שעד כה לא השתמשנו
בו למרות שהוא קריטי: יש לכל המרחבים נקודה משותפת, \L{$a\in\bigcap_{i\in\Lambda}W_{i}$}.
עבור ה-\L{$W_{k}$} שלנו שעבורו \L{$W_{k}=A_{k}\cup B_{k}$} והנחנו
בלי הגבלת הכלליות ש-\L{$B_{k}=\emptyset$} חייב להתקיים \L{$a\in A_{k}\subseteq A$}.
אבל עכשיו, עבור \L{$W_{i}=A_{i}\cup B_{i}$} אחר כלשהו, \textbf{עדיין}
חייב להתקיים \L{$a\in A_{i}$}; כי אם היה מתקיים \L{$a\in B_{i}\subseteq B$}
היינו מקבלים \L{$A\cap B\ne\emptyset$} בסתירה לכך ש-\L{$A,B$} זרות.
אז המסקנה היא שגם כאן, \L{$A_{i}\ne\emptyset$} ו-\L{$B_{i}=\emptyset$}
וזה לכל \L{$i\in\Lambda$}, כך שסיימנו את ההוכחה.

אז בניה אחת עברה בשלום - איחוד של מרחבים קשירים עם נקודה משותפת הוא
קשיר. הדבר המתבקש הבא לטעמי הוא מכפלה, אבל זה יהיה קצת יותר טריקי,
אז בואו נדבר קודם על משהו אחר, מתבקש קצת פחות: אם \L{$A$} היא קבוצה
קשירה, אז גם \L{$\overline{A}$}, הסגור של \L{$A$}, היא קבוצה קשירה.
כזכור, \L{$\overline{A}$} הוא מה שמתקבל מ-\L{$A$} כשמוסיפים את כל
נקודות השפה - אלו שבכל סביבה שלהן יש משהו מתוך \L{$A$} ומשהו מחוץ
ל-\L{$A$}. מכיוון שנקודות השפה של \L{$A$} \textquotedblleft צמודות\textquotedblright{}
ככה ל-\L{$A$} עצמה, הגיוני שלהוסיף אותן לא יאפשר פתאום לפצל את \L{$\overline{A}$}
לשתי קבוצות פתוחות זרות. אפשר, כמובן, לדמיין מצב שבו יש לנו עיגול
פתוח ואנחנו מוסיפים לו את השפה שלו שהיא מעגל, ואז התוצאה בנויה משתי
חתיכות זרות - העיגול, והמעגל שמהווה את השפה שלו. אבל המעגל הזה הוא
לא קבוצה פתוחה - יש בה נקודות שהן לא \textquotedblleft עמוק בפנים\textquotedblright{}
אלא כל סביבה שלהן יוצאת החוצה מהקבוצה. אז הגיון יש פה, אבל איך הפורמליזם
עובד?

ובכן, בואו נניח שיש פירוק \L{$\overline{A}=B\cup C$}. איפה \L{$A$}
נמצאת ביחס לפירוק הזה? בוודאי ש-

\L{$A=\left(B\cup C\right)\cap A=\left(B\cap A\right)\cup\left(C\cap A\right)$}

אבל בגלל ש-\L{$A$} קשירה, אחת משתי הקבוצות הללו חייבת להיות ריקה,
למשל \L{$C\cap A=\emptyset$} ולכן \L{$A\subseteq B$}. עכשיו ניקח
את הסגור של שתי הקבוצות הללו ונקבל \L{$\overline{A}\subseteq\overline{B}$}
)לא הוכחתי שלקחת סגור לשתי הקבוצות משמר את ההכלה, אבל קל לראות את
זה - סגור של קבוצה הוא חיתוך כל הקבוצות הסגורות שמכילות אותה, וכל
מי שמכיל את \L{$B$} מכיל גם את \L{$A$} ולכן החיתוך שמגדיר את \L{$\overline{A}$}
כולל \textbf{לכל הפחות} את אותן קבוצות בחיתוך שמגדיר את \L{$\overline{B}$}(.

עכשיו, התחלנו מזה ש-\L{$\overline{A}=B\cup C$}. בגלל ש-\L{$B,C$}
זרות, זה אומר ש-\L{$B=\overline{A}\backslash C$} )כי כל מה שלא ב-\L{$C$}
חייב להיות ב-\L{$B$}, וההפך(. אז מה קיבלנו פה? תזכרו שהמרחב שבתוכו
אנחנו עובדים כרגע הוא \L{$\overline{A}$} - זה ה\textquotedblright עולם\textquotedblright{}
הטופולוגי שלנו. הקבוצה \L{$C$} פתוחה ביחס לעולם הזה, ולכן אם \L{$B$}
היא המשלימה של \L{$C$} ביחס לעולם הזה, \L{$B$} היא סגורה ולכן שווה
לסגור שלה בעולם הזה. הסגור הזה הוא \L{$\overline{A}\cap\overline{B}$}
)כי הסגור של \L{$B$} ביחס לתת-מרחב מסוים הוא החיתוך של הסגור של \L{$B$}
במרחב כולו עם תת המרחב הזה( ולכן קיבלנו \L{$\overline{A}\subseteq\overline{A}\cap\overline{B}=B$}ולכן
\L{$C$} חייבת להיות ריקה, מה שמסיים את ההוכחה הזו.

אפשר לקחת את הרעיון הזה צעד אחד קדימה - לא רק \L{$\overline{A}$}
תהיה קשירה, אלא כל קבוצה שמתקבלת מהוספת \textquotedblleft חלק מהשפה\textquotedblright{}
של \L{$A$} אל \L{$A$}, כלומר כל \L{$A\subseteq A^{\prime}\subseteq\overline{A}$}.
הרעיון זהה: הפעם \L{$A^{\prime}$} תהיה ה\textquotedblright עולם\textquotedblright{}
שלנו. לוקחים פירוק \L{$A^{\prime}=B\cup C$}, מסיקים ש-\L{$A$} מוכלת
ב-\L{$B$} ולכן \L{$\overline{A}\subseteq\overline{B}$} . כעת \L{$B=A^{\prime}\backslash C$}
סגורה ב-\L{$A^{\prime}$} ולכן שווה לסגור שלה, \L{$A^{\prime}\cap\overline{B}$},
ולכן נסיק \L{$A^{\prime}\subseteq A^{\prime}\cap\overline{B}=B$}
ונסיים באותו האופן.

אוקיי, עכשיו - מה עם מכפלה? קל לדמיין מכפלה בהקשר של \L{$\mathbb{R}$}:
למשל, המכפלה של הקטע \L{$\left[0,1\right]$} בקטע \L{$\left[3,6\right]$}
היא המלבן במישור \L{$\left\{ \left(x,y\right)\ |\ 0\le x\le1,3\le y\le6\right\} $}.
די ברור שאם נכפול משהו לא קשיר במשהו קשיר נקבל משהו לא קשיר, למשל
איחוד של שני קטעים זרים שמוכפל בקטע אחד יתן לנו שני מלבנים זרים. מצד
שני, מכפלה של שתי קבוצות קשירות - כלומר, של שני קטעים - כן נותנת לנו
קבוצה קשירה. אז מתבקש להניח שזה קורה באופן כללי, שמכפלה של מרחבים
קשירים היא קשירה. זה למעשה נכון, אבל נשאלת השאלה - \textbf{איזו} מכפלה?
או ליתר דיוק - איזו \textbf{טופולוגיה} יש על מרחב המכפלה?

\href{https://gadial.net/2026/07/25/topological_spaces_subspace_and_product_topology/}{כשהצגתי מרחבי מכפלה},
היו שתי טופולוגיות שאפשר להגדיר באופן \textquotedblleft טבעי\textquotedblright{}
על מרחבים כאלו: טופולוגיית התיבה וטופולוגיית המכפלה. מבין שניהם טופולוגיית
התיבה נראתה הכי טבעית: קבוצה פתוחה היא פשוט מכפלה של קבוצות פתוחות,
וזהו. טופולוגיית המכפלה הייתה קצת יותר מוזרה - קבוצה פתוחה היא מכפלה
של קבוצות פתוחות אבל כל הקבוצות הללו פרט למספר סופי חייבות להיות \textbf{כל
המרחב}. היה איזה הסבר טוב להגדרה הזו - משהו עם תת-בסיסים והטלות -
אבל נעזוב את זה; הפואנטה הייתה שבטופולוגיית המכפלה קורים דברים יפים,
כמו משפט טיכונוף שאומר שמכפלה של מרחבים קומפקטיים היא קומפקטית, ובטופולוגיית
התיבה זה לא קורה. זה בדיוק המצב גם עם קשירות: בטופולוגיית התיבה, מכפלה
של מרחבים קשירים לא חייבת להיות קשירה, ובטופולוגיית המכפלה כן.

בואו נתחיל עם הדוגמא הנגדית עבור טופולוגיית התיבה. מן הסתם, ממש פשוט
זה לא יהיה - זה חייב לערב מכפלה אינסופית, אבל בואו נשמור את זה פשוט
ככל הניתן ונכפול מספר בן מניה של עותקים של \L{$\mathbb{R}$}, כלומר
נסתכל על הקבוצה \L{$\mathbb{R}^{\omega}$} שאפשר לחשוב על האיברים
שלה פשוט בתור סדרות אינסופיות \L{$\left\{ a_{n}\right\} ^{\infty}_{n=0}$}
של מספרים ממשיים. עכשיו, סדרה של ממשיים יכולה להיות אחד משניים: או
שהיא \textbf{חסומה}, כלומר קיים \L{$M$} כך ש-\L{$\left|a_{i}\right|\le M$}
לכל \L{$i\in\mathbb{N}$}, או שהיא לא. נסמן ב-\L{$A$} את קבוצת כל
הסדרות החסומות וב-\L{$B$} את קבוצת כל הסדרות הלא חסומות. מן הסתם
הן זרות ומשלימות - האיחוד שלהן נותן את כל \L{$\mathbb{R}^{\omega}$},
ולכן אם \L{$A,B$} שתיהן פתוחות אנחנו בצרה צרורה כי מצאנו פירוק של
המרחב ומכאן שהוא לא קשיר.

כדי להראות שקבוצה היא פתוחה, ניקח סדרה \L{$a_{0},a_{1},a_{2},\ldots$}
כללית ונמצא קבוצה פתוחה שמכילה אותה ומוכלת כולה באחת מהקבוצות. אנחנו
לא צריכים אפילו להתאמץ עם אפסילונים הפעם, פשוט נסתכל על הקבוצה

\L{$\left(a_{0}-1,a_{0}+1\right)\times\left(a_{1}-1,a_{1}+1\right)\times\ldots$}

שכוללת את כל הסדרות שהאיברים שלהן, בכל קואורדינטה, הם לכל היותר במרחק
{\beginL 1\endL} )בחירה שרירותית לגמרי( מאברי הסדרה המקורית. זו קבוצה
פתוחה כי היא מכפלה של קבוצות פתוחות ב-\L{$\mathbb{R}$}. עכשיו, מן
הסתם אם הסדרה המקורית הייתה חסומה, גם כל סדרה שמתקבלת ממנה בכך שזזים
ב-{\beginL 1\endL} לכל היותר בכל מקום עדיין תהיה חסומה, ואם הסדרה
המקורית הייתה לא חסומה גם כל סדרה שמתקבלת ממנה ככה לא תהיה חסומה -
אז הקבוצה הזו מוכלת כולה או ב-\L{$A$} או ב-\L{$B$}, בהתאם לשאלה
באיזה משתי הקבוצות הללו הסדרה המקורית שלנו נמצאת. זה מסיים את הסיפור
עבור טופולוגיית התיבה, אם כי ברור שהבניה הזו לא תעבוד לטופולוגיית
המכפלה כי כמעט כל המוכפלים יהיו חייבים להיות \L{$\mathbb{R}$}, כלומר
בכל קבוצה פתוחה שמכילה את הסדרה שלנו יוכלו להיות גם סדרות חסומות וגם
סדרות לא חסומות.

בואו נעבור עכשיו להוכחה עבור המקרה שבו זה עובד. ראשית, מקרה הבסיס
הפשוט שלנו יהיה מכפלה \L{$A\times B$} של שני מרחבים קשירים. אם נוכיח
את זה, אפשר לקבל באינדוקציה שמכפלה של מספר \textbf{סופי} כלשהו של
מרחבים קשירים הוא קשיר )כזכור, במקרה הסופי טופולוגיית התיבה וטופולוגיית
המכפלה נותנות את אותה הטופולוגיה(. כאן אנחנו קוטפים את הפירות של מה
שכבר ראינו - שאיחוד של מרחבים קשירים עם נקודה משותפת הוא קשיר, ושתמונה
רציפה של מרחב קשיר היא קשירה.

ראשית, בואו נבחר נקודה \L{$\left(a,b\right)\in A\times B$} כלשהי
)אם אין, כי \L{$A$} או \L{$B$} ריקה, זה מקרה מנוון אבל לא בעייתי
כי טכנית הפירוק \L{$\emptyset=\emptyset\cup\emptyset$} מראה שהקבוצה
הריקה קשירה(. עכשיו, בואו נדמיין שניה ש-\L{$A$} ו-\L{$B$} שניהם
\L{$\mathbb{R}$} ולכן אפשר לחשוב עליהם בתור קווים ישרים ועל המכפלה
שלהם בתור המישור ולכן על \L{$\left(a,b\right)$} בתור נקודה במישור.
מה זו הקבוצה \L{$A\times b$}? זה הקו \L{$A$} )שנדמיין בתור קו אופקי(
שעובר דרך ראשית הצירים בקואורדינטה \L{$\left(0,b\right)$}. באופן
דומה \L{$a\times B$} הוא קו אנכי, ולכן האיחוד של שניהם \L{$\left(A\times b\right)\cup\left(a\times B\right)$}
הוא מין \textquotedblleft צלב\textquotedblright{} כזה שבו נמצאת הנקודה
\L{$\left(a,b\right)$} שמשותפת לשני המרחבים. עכשיו, \L{$A\times b$}
הומיאומורפי ל-\L{$A$} )על ידי הפונקציה \L{$f\left(x\right)=\left(x,b\right)$}
שקל להראות שהיא חח\textquotedblright ע, על ורציפה( ולכן אם \L{$A$}
קשיר גם \L{$A\times b$} קשיר; מאותו נימוק \L{$a\times B$} קשיר;
ולכן ה\textquotedblright צלב\textquotedblright{} \L{$\left(A\times b\right)\cup\left(a\times B\right)$}
הוא קשיר בעצמו. רק מה, הצלב הזה לא מכסה את כל \L{$A\times B$}, צריך
עוד קבוצות. אבל זו התחלה טובה.

כדי להמשיך, בואו נבחר \L{$a^{\prime}\in A$} כלשהו ונגדיר צלב דומה,
\L{$\left(A\times b\right)\cup\left(a^{\prime}\times B\right)$}.
גם הוא יהיה קשיר מאותו נימוק כמו קודם. עכשיו בואו נאחד את כל הצלבים
שנוצרים באופן הזה: \L{$\bigcup_{a^{\prime}\in A}\left(A\times b\right)\cup\left(a^{\prime}\times B\right)$}.
מה קיבלנו? ראשית, קיבלנו את כל המרחב; עבור נקודה שרירותית כלשהי \L{$\left(x,y\right)$}
היא שייכת אל \L{$a^{\prime}\times B$} עבור המקרה \L{$x=a^{\prime}$}.
שנית, כל הצלבים חולקים נקודה משותפת כי כולם חולקים את \L{$A\times b$}
- תדמיינו את זה כאילו יש לנו ישר אופקי קבוע, וישר אנכי שעובר דרכו
וזז ימינה ושמאלה עד שהוא מכסה את כל המישור. זה מסיים את ההוכחה במקרה
של מכפלה של שני מרחבים, ולכן של מספר סופי כלשהו.

נשאר רק המקרה של מכפלה אינסופית עם טופולוגיית המכפלה. הרעיון הבסיסי
מאוד דומה למה שראינו במקרה הסופי - אנחנו עדיין עובדים עם \textquotedblleft צלבים\textquotedblright{}
קשירים שכאלו עם נקודה אחת משותפת ומנסים לתפוס באמצעות איחוד שלהם כמה
שיותר מהמרחב. ההבדל הוא שהפעם האיחוד לא ייתן לנו את \textbf{כל} המרחב,
אבל הוא ייתן לנו קבוצה ש\textbf{הסגור} שלה הוא כל המרחב, וכבר ראינו
שאם קבוצה היא קשירה גם הסגור שלה קשיר.

אם כן, בואו ניקח מכפלה כללית \L{$\prod_{i\in\Lambda}A_{i}$} של קבוצות
קשירות. אם היא ריקה, קיבלנו מרחב קשיר וסיימנו. אחרת, יש בה נקודה ספציפית
\L{$a=\left\{ a_{i}\right\} _{i\in\Lambda}$}. עכשיו, בואו ניקח תת-קבוצה
\textbf{סופית} של האינדקסים, \L{$F\subseteq\Lambda$}, ונסתכל על \textquotedblleft צלב\textquotedblright{}
שנבנה בעזרת \L{$F$}: הקבוצה 

\L{$C_{F}=\left\{ x\in\prod_{i\in\Lambda}A_{i}\ |\ i\notin F\Rightarrow x_{i}=a_{i}\right\} $}

כלומר, זו קבוצת כל הנקודות במכפלה ש\textbf{זהות} ל-\L{$a$} בכל האינדקסים
למעט אולי הכמות הסופית של האינדקסים שיש ב-\L{$F$}. לא קשה לראות שהקבוצה
הזו הומיאומורפית למכפלה \textbf{הסופית} \L{$\prod_{i\in F}A_{i}$}
והמכפלה הזו קשירה על פי מה שכבר ראינו. בנוסף, \L{$a\in C_{F}$}. לכן,
אם ניקח את \textbf{כל} ה-\L{$C_{F}$}-ים האפשריים שמתקבלים מלקיחת
קבוצה סופית \textbf{כלשהי} \L{$F\subseteq\Lambda$}, נקבל איחוד של
מרחבים קשירים עם נקודה משותפת )\L{$a$}(. זה איחוד עצום, שיכול להיות
גם של מספר לא בן מניה של מרחבים; זה לא משנה, ההוכחה שלנו על איחוד
עבדה גם במקרים הללו. אז אפשר להגדיר \L{$C=\bigcup_{F}C_{F}$} ולקבל
תת-מרחב קשיר \L{$C\subseteq\prod_{i\in\Lambda}A_{i}$}. נשאר רק להראות
ש-\L{$\overline{C}$} כולל את כל \L{$\prod_{i\in\Lambda}A_{i}$}.
זה דורש להראות שלכל איבר במכפלה יש סביבה פתוחה עם חיתוך לא ריק עם
\L{$C$}, אבל אפשר למעשה להראות יותר מזה: ש-\L{$C$} \textbf{צפופה}
במכפלה. אני לא חושב שהזכרתי את המושג הזה עדיין, אבל הוא פשוט: \L{$C$}
צפופה במרחב כלשהו אם כל קבוצה פתוחה לא ריקה במרחב היא בעלת חיתוך לא
ריק עם \L{$C$}. הדוגמא הקלאסית היא \L{$\mathbb{Q}$} שצפופה במרחב
\L{$\mathbb{R}$}.

אם כן, בואו ניקח קבוצה פתוחה לא ריקה כללית \L{$U=\prod_{i\in\Lambda}U_{i}$}
במרחב המכפלה שלנו. כאן נכנסת לתמונה העובדה שאנחנו עובדים בטופולוגיית
המכפלה - הקבוצה הפתוחה הזו חייבת להיות \textbf{די פשוטה}, כלומר \L{$U_{i}=A_{i}$}
לכל האינדקסים חוץ ממספר סופי. רואות לאן זה הולך? נסמן את קבוצת האינדקסים
שבהם \L{$U_{i}\ne A_{i}$} ב-\L{$F$}, ועכשיו ניקח איבר \L{$x\in\prod_{i\in\Lambda}U_{i}$}
כלשהו ונתעלל באינדקסים שלו - נשנה את כולם לאינדקסים של \L{$a$} חוץ
מאשר את האינדקסים ב-\L{$F$}. כלומר, נגדיר

\L{$x^{\prime}_{i}=\begin{cases}
x_{i} & i\in F\\
a_{i} & i\notin F
\end{cases}$}

ועכשיו \L{$x^{\prime}\in C_{F}$} ממש על פי ההגדרה. בנוסף, לא ברחנו
מגבולות \L{$U$}; עדיין מתקיים \L{$x^{\prime}\in U$} כי בכל אינדקס
שאינו ב-\L{$F$}, הערך של אברי \L{$U$} \textbf{הוא בלתי מוגבל}. ההוכחה
הזו נחמדה כי אנחנו רואים בדיוק את האופן שבו טופולוגיית המכפלה מאפשרת
לנו \textquotedblleft לקפוץ\textquotedblright{} מהמקרה הסופי למקרה
הכללי: במקרה הסופי הצלחנו להוכיח משהו לכל קבוצה סופית \L{$F$}, וטופולוגיית
המכפלה מאפשרת לנו להסתפק \textbf{רק} בכך שהוכחנו עבור \L{$F$} סופית
ולא עבור דברים גדולים יותר.

\section*{קשירות מסילתית}

אם כבר אני כותב פוסט על קשירות, בואו ננצל את שולי הפוסט כדי להזכיר
מושג קשור ורלוונטי בפני עצמו, שאולי אפילו יותר מתאים לאינטואיציה הראשונית
של \textquotedblleft קשירות\textquotedblright{} - \textbf{קשירות מסילתית}.
בשביל זה אולי נוח לדמין גרף - בתורת הגרפים, גרף הוא \textbf{קשיר}
אם יש מסלול בין כל שני צמתים בו. אם אנחנו מציירים גרף עם צמתים בתור
נקודות וקשתות בתור קווים, אז מסלול בגרף הוא פשוט קו שבור ארוך, ו\textquotedblright קשירות\textquotedblright{}
מתבטאת בכך שיש קו כזה שעובר כולו בתוך הגרף עבור כל שתי נקודות בגרף.

קשירות מסילתית היא רעיון דומה - מרחב \L{$X$} הוא קשיר מסילתית אם
לכל זוג נקודות \L{$x,y\in X$} יש מסלול רציף שכזה שמגיע מ-\L{$x$}
אל \L{$y$}. כמובן, יש פה בעיה קטנה - מרחב טופולוגי כללי \L{$X$}
הוא משהו אבסטרקטי שכזה, איך נראה \textquotedblleft מסלול\textquotedblright{}
בו? והתשובה לזה פשוטה ודומה להרבה תשובות אחרות; בואו נסתכל על \L{$\mathbb{R}$}
שבו אנחנו כן מבינים דברים ואז נזרום משם הלאה עם פונקציות רציפות.

ב-\L{$\mathbb{R}$} מסלול בין שתי נקודות \L{$a<b$} הוא פשוט הקטע
\L{$\left[a,b\right]$} שמתחיל ב-\L{$a$} ונגמר ב-\L{$b$} ועובר בכל
הנקודות באמצע. כמובן, אפשר לדמיין מסלולים מסובכים יותר שהולכים קדימה
ואחורה ונעצרים קצת באמצע לראות את הנוף וכדומה, אבל גם אותם אפשר לקבל
מהמסלול הפשוט \L{$\left[a,b\right]$}. פורמלית: במרחב \L{$X$} נגדיר
מסלול מ-\L{$x$} אל \L{$y$} בתור פונקציה רציפה \L{$\gamma:\left[a,b\right]\to X$}
עבור \L{$a,b\in\mathbb{R}$}, כך ש-\L{$\gamma\left(a\right)=x$} ו-\L{$\gamma\left(b\right)=y$}.

בואו נראה דוגמא פשוטה. ניקח \L{$X=\mathbb{R}^{2}$} כדי שעדיין יהיה
משהו שאפשר לדמיין, ונבחר את הנקודות שעל הצירים, \L{$x=\left(1,0\right)$}
ו-\L{$y=\left(0,1\right)$}. עכשיו בואו ניקח מעגל היחידה שעובר דרך
שתי הנקודות הללו - אם נסתכל רק על הרבע שלו שהולך מ-\L{$x$} אל \L{$y$},
יש לנו דוגמא למסלול בין שתי הנקודות הללו. איך נגדיר אותו פורמלית?
ניקח את הקטע \L{$\left[0,\frac{\pi}{2}\right]$} ונגדיר עליו \L{$\gamma\left(t\right)=\left(\cos t,\sin t\right)$}
- זו הפרמטריזציה הסטנדרטית של מעגל, כך שאין פה משהו מרגש - אבל זה
עובד יפה.

עכשיו בואו נראה שאם \L{$X$} קשיר מסילתית אז הוא קשיר. ניקח פירוק
כלשהו, \L{$X=A\cup B$} כך ש-\L{$A,B$} קבוצות פתוחות זרות. נניח בשלילה
ששתיהן לא ריקות, כך שיש \L{$x\in A,y\in B$}. מכיוון שהמרחב קשיר מסילתית,
קיימת פונקציה רציפה \L{$\gamma:\left[a,b\right]\to X$} כך ש-\L{$\gamma\left(a\right)\in A$}
ו-\L{$\gamma\left(b\right)\in B$}. מצד שני, \L{$\gamma$} רציפה והקטע
\L{$\left[a,b\right]$} קשיר )לא הוכחתי את זה בפוסט הזה! אני מסתמך
על זה שאנחנו כבר יודעים שאקסיומת השלמות נותנת את זה( ולכן התמונה \L{$\gamma\left(\left[a,b\right]\right)$}
היא תת-קבוצה קשירה של \L{$X$}. כבר ראינו בהוכחות קודמות שזה אומר
שהיא מוכלת כולה ב-\L{$A$} או ב-\L{$B$} )אחרת החיתוך שלה עם שניהם
יתן פירוק שלה( וזה בסתירה לכך ש-\L{$\gamma\left(a\right)\in A$} ו-\L{$\gamma\left(b\right)\in B$}.
הוכחה פשוטה! וגם די אינטואיטיבית, הרי אם המרחב מפורק לשתי חתיכות לא
נוכל למצוא מסלול רציף מאחת לשניה.

מה שמעניין פה הוא שההפך לא נכון. מרחב יכול להיות קשיר אבל לא קשיר
מסילתית, כי, ובכן, טופולוגיה מתעסקת לפעמים עם מרחבים מוזרים שבהם האינטואיציה
שלנו לא עובדת. זה ממחיש למה ההגדרה הטופולוגית לקשירות, עם הפירוק הזה
לשתי קבוצות, יותר רלוונטית )לטופולוגיה( מאשר זו של קשירות מסילתית
שנראית אולי טבעית יותר )והיא שימושית מאוד בפני עצמה, לא סתם טרחו לתת
לה שם(. אלא שכדי להיות אקסטרה מעליבים, הדוגמא הנגדית היא אפילו לא
כזו מוזרה - היא עקום די פשוט שחי לו בסבבה ב-\L{$\mathbb{R}^{2}$}
)לא צריך מרחבים שהם מכפלות אינסופיות מוזרות( ונקרא בשם היפה והקולע
\textquotedblleft עקומת הסינוס של הטופולוגים\textquotedblright . ככה
זה נראה:

\textbf{איור}

הדבר הזה הוא בסך הכל הגרף של הפונקציה \L{$f\left(x\right)=\sin\frac{1}{x}$}
בקטע \L{$(0,1]$}, כלומר אוסף הנקודות מהצורה \L{$\left(x,\sin\frac{1}{x}\right)$}
עבור \L{$0<x\le1$}, יחד עם הקו האנכי שעובר דרך ראשית הצירים, \L{$\left\{ 0\right\} \times\left[-1,1\right]$},
שאותו מוסיפים כדי לקבל את הסגור של הגרף של הסינוס. תחשבו על זה ככה:
ככל ש-\L{$x$} מתקרב ל-\L{$0$}, \L{$\frac{1}{x}$} שואף לאינסוף יותר
ויותר מהר, ולכן \L{$\sin\frac{1}{x}$}, שתמיד \textquotedblleft מזפזפ\textquotedblright{}
בין כל הערכים האפשריים בין \L{$1$} ובין \L{$-1$} ממשיך לזפזפ בינם
יותר ויותר מהר, ולכן לכל \L{$\varepsilon>0$} קטן כרצוננו, תמיד אפשר
יהיה למצוא נקודה על העקום שעבורה \L{$x<\varepsilon$} ואילו \L{$\sin\frac{1}{x}$}
הוא נקודה \L{$t\in\left[-1,1\right]$} כלשהי ולכן המרחק של הנקודה
על העקום מהנקודה \L{$\left(0,t\right)$} על הקו האנכי הוא קטן מ-\L{$\varepsilon$},
כלומר \L{$\left(0,t\right)$} היא נקודת הצטברות של גרף הסינוס, ולכן
אם לוקחים את הסגור שלו צריך להוסיף את כל הנקודות מהצורה הזו. אני אסמן
\L{$S=\left\{ \left(x,\sin\frac{1}{x}\right)\ |\ 0<x\le1\right\} $}
ואז עקומית הסינוס של הטופולוגים היא \L{$\overline{S}$}.

עכשיו, זה ש-\L{$\overline{S}$} היא קבוצה קשירה נובע מיידית מדברים
שכבר ראינו: הגרף של \L{$\sin\frac{1}{x}$} הוא תמונה רציפה של הקטע
הקשיר \L{$(0,1]$} ולכן קשיר; וכשלוקחים את הסגור שלו, מקבלים קבוצה
קשירה. אבל מצד שני, הגרף של הסינוס לא באמת \textbf{נוגע} בקו האנכי,
נכון? הוא רק מתקרב אליו באופן מעוך שכזה. זה מספיק לקשירות אבל זה \textbf{לא}
מספיק לקשירות מסילתית. עדיין, איך מראים את זה פורמלית?

בואו נניח שהמרחב כן קשיר מסילתית וניקח מסלול \L{$\gamma:\left[a,b\right]\to\overline{S}$}
כך ש-\L{$\gamma\left(a\right)\in\overline{S}\backslash S$} )כלומר,
בקו האנכי( ו-\L{$\gamma\left(b\right)\in S$} )כלומר, על גרף הסינוס
עצמו(. איך המסלול נראה? בהתחלה הוא יכול לזוז אנה ואנה בקו האנכי, כמו
מישהו שמתרוצץ על שפת נהר. אולי הוא אפילו יעבור באורח קסום אל \L{$S$}
ואז יחזור אל הקו האנכי, מי יודע? מה שאנחנו רוצים לעשות הוא להתמקד
בשלב האחרון - הרגע שבו הוא עוזב את הקו האנכי על מנת שלא לחזור עוד.
המושג הזה של \textquotedblleft רגע\textquotedblright{} מתאים יפה לחשיבה
על \L{$\gamma:\left[a,b\right]\to\overline{S}$} בתור פונקציה שממפה
\textbf{זמן} למקום, כלומר אנחנו חושבים על \L{$\gamma\left(t\right)$}
בתור \textquotedblleft איפה המסלול נמצא בזמן \L{$t$}\textquotedblright .
אז בואו נסתכל על קבוצת \textquotedblleft כל הרגעים שבהם המסלול נמצא
בקו האנכי\textquotedblright : זה בדיוק \textbf{המקור} של הקו האנכי
על פי \L{$\gamma$}, כלומר הקבוצה \L{$\gamma^{-1}\left(\left\{ 0\right\} \times\left[-1,1\right]\right)$}.
עכשיו, מכיוון שהקו האנכי הוא קבוצה סגורה ו-\L{$\gamma$} רציפה, גם
המקור שלה הוא קבוצה סגורה. מכיוון שהיא חסומה )על ידי \L{$b$}(, יש
לה \textbf{חסם עליון}, ומכיוון שהיא סגור החסם העליון הזה שייך לקבוצה
עצמה, כלומר גם הוא איבר שמועבר על ידי \L{$\gamma$} אל הקו האנכי.
נסמן אותו ב-\L{$c$}.

מכיוון ש-\L{$\gamma\left(b\right)\in S$} ו-\L{$\gamma\left(c\right)\in\overline{S}\backslash S$},
אנחנו יודעים ש-\L{$c<b$}, מה שאומר שאפשר עכשיו להסתכל רק על \L{$\gamma:\left[c,b\right]\to\overline{S}$}:
זה מסלול שברגע הראשון שלו הוא לא ב-\L{$\overline{S}$} אבל בכל רגע
אחר הוא כן.

אבל מה הבעיה פה? מן הסתם, משהו ברציפות של \L{$\gamma$} חייב להשתבש.
מכיוון ש-\L{$\gamma$} היא פונקציה עם שני רכיבים, אפשר לכתוב אותם
במפורש: \L{$\gamma\left(t\right)=\left(x\left(t\right),y\left(t\right)\right)$}.
מספיק יהיה להראות שהרכיב \L{$y$} הוא לא רציף, אבל גם ב-\L{$x$} יהיה
לנו שימוש.

דרך פשוטה לשבש רציפות של פונקציה הוא למצוא \textbf{סדרה} של נקודות
שמתכנסות לנקודה שבה אנחנו רוצים לרסק את הרציפות, וסדרת \textbf{הפלטים}
על הנקודות הללו בבירור לא מתכנסת אל הערך של הפונקציה בנקודה הזו. אצלנו
הנקודה הבעייתית היא \L{$c$}: נרצה לבנות סדרה \L{$\left\{ t_{n}\right\} ^{\infty}_{n=1}$}
כך ש-\L{$\lim_{n\to\infty}y\left(t_{n}\right)\ne y\left(c\right)$}
למרות ש-\L{$\lim_{n\to\infty}t_{n}=c$}. עם סינוס שמתפרע זה קל במיוחד:
אני אבנה את הסדרה שלי כך ש-\L{$y\left(t_{n}\right)=\left(-1\right)^{n}$}
)כלומר, או {\beginL 1\endL} או \L{$-1$}, בהתאם לזוגיות של \L{$n$}(
ואקבל סדרה שפשוט \textbf{לא מתכנסת}.

איך נעשה את זה? כמו שכבר אמרתי, הקטע עם \L{$\sin\frac{1}{x}$} הוא
שככל שמתקרבים ל-{\beginL 0\endL} הוא מזפזפ מהר יותר ויותר. אבל הנה
הפורמליזם: לכל \L{$n$}, בואו נסתכל על \L{$\left[0,x\left(c+\frac{1}{n}\right)\right]$},
כלומר על כל הנקודות על ציר \L{$x$} החל מ-{\beginL 0\endL} ועד לנקודה
שאליה המסלול מגיע בזמן \L{$c+\frac{1}{n}$}. שימו לב שהקטע הזה הוא
לא טריוויאלית \L{$\left\{ 0\right\} $} ותו לא, בגלל שלכל נקודה גדולה
מ-\L{$c$}, ה-\L{$x$} על הנקודה הזו הוא לא {\beginL 0\endL} )כי הפלט
של \L{$\gamma$} על \L{$c+\frac{1}{n}$} לא נמצא על הקו האנכי(. לכן,
מהתכונה של \L{$\sin\frac{1}{x}$} שתיארתי, קיימת בקטע הזה נקודה \L{$s\in\left[0,x\left(c+\frac{1}{n}\right)\right]$}
כך ש-\L{$\sin\frac{1}{s}=\left(-1\right)^{n}$}. מה נשאר? למצוא \L{$c<t_{n}\le c+\frac{1}{n}$}
כך ש-\L{$x\left(t_{n}\right)=s$}. למה זה קיים? ובכן, נכון ש-\L{$x$}
היא רכיב של \L{$\gamma$} הרציפה ולכן היא בעצמה פונקציה רציפה? ומתקיים
\L{$x\left(c\right)=0$}? ומתקיים, ובכן, \L{$x\left(c+\frac{1}{n}\right)=x\left(c+\frac{1}{n}\right)$}?
אז יש לנו פונקציה \textbf{רציפה} שמוגדרת על קטע סגור, ויש לנו \textbf{ערך
ביניים }\L{$s\in\left[x\left(c\right),x\left(c+\frac{1}{n}\right)\right]$},
ולכן אנחנו יודעים שקיים \L{$t_{n}$} כך ש-\L{$x\left(t_{n}\right)=s$},
ממשפט ערך הביניים שאיתו התחלנו את כל הסיפור. אמרתי שהוא הולך להיות
שימושי!
\end{document}
